\documentclass[12pt]{article}
\usepackage[paperheight=210mm,paperwidth=148mm,top=2cm,bottom=2cm,left=1cm,right=1cm]{geometry}  % A5 paper
\usepackage{background}
\backgroundsetup{
scale=1,
opacity=0.1,
angle=0,
color=black,
contents={%
\includegraphics[width=\paperwidth]{background}
}
}
%%%%%%%%%%%%%%%%%%%%%%%%%%%%%%%%%%%%%%%%
\usepackage[utf8]{inputenc}
\usepackage{titlesec}
\usepackage{multirow}
\usepackage{tikz}
\usepackage{ctable}
\usepackage{booktabs}
%%%%%%%%%%%%%%%%%%%%%%%%%%%%%%%%%%%%%%%%
\renewcommand{\familydefault}{\sfdefault}
% \titleformat{\chapter}[display]
% {\bfseries\LARGE\sffamily}{\chaptertitlename}{10pt}{\LARGE}
% \titlespacing*{\chapter}{0pt}{10pt}{10pt}
% \numberwithin{table}{chapter}


\begin{document}
% Cover
\newgeometry{left=0cm, bottom=0cm, top=0cm, right=0cm}
\thispagestyle{empty}
\noindent  % To remove the unwanted white space.
\includegraphics{capa}
\NoBgThispage

\clearpage

\restoregeometry

\newpage

% Map
\newgeometry{left=0cm, bottom=0cm, top=0cm, right=0cm}
\thispagestyle{empty}
\noindent  % To remove the unwanted white space.
\includegraphics[width=\textwidth]{../../media/img/mapa_UFSC_2014_1500x2121}

\clearpage

\restoregeometry

\section*{Letter from organizers}

Hi,


\newpage

\section*{Meet our Sponsors}
\subsection*{Silver Sponsors}
\begin{minipage}{0.45\textwidth}
\textbf{Python Software Foundation}

\includegraphics[width=\textwidth]{../../media/img/psf}

The Python Software Foundation (PSF) is a 501(c)(3) non-profit corporation that
holds the intellectual property rights behind the Python programming language.
We manage the open source licensing for Python version 2.1 and later and own and
protect the trademarks associated with Python. We also run the North American
PyCon conference annually, support other Python conferences around the world,
and fund Python related development with our grants program and by funding
special projects.

\end{minipage}
\hfill
\begin{minipage}{0.45\textwidth}
\textbf{NumFOCUS}

\includegraphics[width=\textwidth]{../../media/img/numfocus}

NumFOCUS is a 501(c)(3) nonprofit that supports and promotes world-class,
innovative, open source scientific software.

The mission of NumFOCUS is to promote sustainable high-level programming
languages, open code development, and reproducible scientific research. We
accomplish this mission through our educational programs and events as well as
through fiscal sponsorship of open source data science projects. We aim to
increase collaboration and communication within the scientific computing
community.

\end{minipage}

\newpage

\subsection*{16/05/2016}

\begin{center}
       \begin{tabular}{@{}l p{5cm} p{5cm}@{}}
     \toprule
      & 9h-12h & 14h-17h\\\midrule
     Room 1 & \textbf{Python for geoscientist} (Filipe Fernandes) & \textbf{Utilización de Vispy para visualización rápida en 3D} (David Ochoa)\\
     Room 2 & \textbf{Introduction to MPI on python using mpi4py} (Jeudy Blanco) & \textbf{Introduction to High-Performance Scientific Computing with PyOpenCL} (Celia Cintas)\\
     Auditorium & \textbf{I see 42 everywhere} (Fernando Masanori) & \textbf{The stupid content tracker} (Tomas Aliaga)\\
     CIASC & \textbf{Software Carpentry Workshop} & \textbf{Software Carpentry Workshop} \\\bottomrule
   \end{tabular}
\end{center}

\subsection*{17/05/2016}

\begin{center}
     \begin{tabular}{@{}l p{5cm} p{5cm}@{}}
     \toprule
      & 9h-12h & 14h-17h\\\midrule
     Room 1 & \textbf{Fit and predict your data: Introduction to Scikit-learn}
     (Celia Cintas) & \textbf{Automatizando a criação de PDF's com LaTeX+python} (Melissa Weber Mendonça)\\
     Room 2 & \textbf{Scientific Python applied to Petroleum Engineering - an
     Introduction} (Igor T. Ghisi) & \textbf{SQL-Alchemy} (Juan Cabral)\\
     Auditorium & \textbf{DSL on the top of Python Object Model} (Juan Cabral) & \textbf{Usando R com o Python - rpy2} (Arnaldo Russo)\\
     CIASC & \textbf{Software Carpentry Workshop} & \textbf{Software Carpentry Workshop} \\\bottomrule
   \end{tabular}
\end{center}

\newpage

\subsection*{18/05/2016}

\begin{center}
   \addtolength{\tabcolsep}{15pt}
   \begin{tabular}{@{}c m{9cm}@{}}
     \toprule
     9:00 \textendash\ 9:40\  & Arrival and registration for lightning talks\\\midrule
     9:40 \textendash\ 10:00 & Welcome\\\midrule
     10:00 \textendash\ 10:40 & Lightning talks\\\midrule
     10:40 \textendash\ 11:20 & Break\\\midrule
     11:20 \textendash\ 12:00 & \textbf{Sinestesia científica - Explorando dados através da música, com DataSounds} (Arnaldo Russo)\\\midrule
     12:00 \textendash\ 12:40 & \textbf{The numba library} (Edison Gustavo Muenz)\\\midrule
     12:40 \textendash\ 14:20 & Lunch Break\\\midrule
     14:20 \textendash\ 15:00 & \textbf{What's new in Spyder 3.0} (Carlos Cordoba)\\\midrule
     15:00 \textendash\ 15:40 & \textbf{External Commands} (Juan Cabral)\\\midrule
     15:40 \textendash\ 16:20 & Break\\\midrule
     16:20 \textendash\ 17:00 & \textbf{import community} (Fernando Masanori)\\\midrule	
     17:00 \textendash\ 19:00 & Happy Hour\\\midrule
     20:00 \textendash\ 23:00 & Dinner\\
     \bottomrule
   \end{tabular}
\end{center}

\clearpage

\subsection*{19/05/2016}

\begin{center}
   \addtolength{\tabcolsep}{15pt}
   \begin{tabular}{@{}c m{9cm}@{}}
     \toprule
     9:00 \textendash\ 9:40 & \textbf{Conda-forge and the future of the Scientific Python packaging} (Filipe Fernandes)\\\midrule
     9:40 \textendash\ 10:20 & \textbf{Geração da Fármacos com Python} (Mário Sérgio de Queiroz)\\\midrule
     10:20 \textendash\ 10:40 & Lightning talks \\\midrule
     10:40 \textendash\ 11:20 & Break \\\midrule
     11:20 \textendash\ 12:00 & \textbf{Fighting techniques in kaggle war} (Emilio Simoni)\\\midrule
     12:00 \textendash\ 12:40 & \textbf{3DBMO: A Python Tool to Study the Power Spectrum Dimensional Dependence in Complex Physical Systems} (Paulo Giovani)\\\midrule
     12:40 \textendash\ 14:20 & Lunch Break \\\midrule
     14:20 \textendash\ 15:00 & \textbf{A Simplificação Textual como instrumento de inclusão cognitiva} (Vinicius Woloszyn)\\\midrule
     15:00 \textendash\ 15:40 & \textbf{Solving Math Programming Problems with Python} (Eduardo Silva)\\\midrule
     15:40 \textendash\ 16:20 & Break \\\midrule
     16:20 \textendash\ 17:00 & Round table about how to run your own event\\\midrule
     17:00 \textendash\ 19:00 & Happy Hour \\\midrule
     20:00 \textendash\ 23:00 & Dinner \\
     \bottomrule
   \end{tabular}
\end{center}

\clearpage

\subsection*{20/05/2016}

\begin{center}
   \addtolength{\tabcolsep}{15pt}
   \begin{tabular}{@{}c m{9cm}@{}}
     \toprule
     9:00 \textendash\ 9:40 & Round table about SciPy Latin America 2017\\\midrule
     9:40 \textendash\ 10:20 & \textbf{Interactive Data Visualization in the Browser with Bokeh} (Bryan Van de Ven)\\\midrule
     10:20 \textendash\ 10:40 & Lightning talks\\\midrule
     10:40 \textendash\ 11:20 & Break\\\midrule
     11:20 \textendash\ 12:00 & \textbf{Ecuaciones diferenciales con Python} (Raul Lopez Briega)\\\midrule
     12:00 \textendash\ 12:40 & \textbf{Localiza, una biblioteca para la localización óptima de equipamientos sociales} (Francisco Palm)\\\midrule
     12:40 \textendash\ 14:20 & Lunch Break\\\midrule
     14:20 \textendash\ 15:00 & \textbf{Learning Analytics: utilizando Data Science para melhorar a educação} (Matheus Braun Magrin)\\\midrule
     15:00 \textendash\ 15:40 & \textbf{From SCADA to IoT with Python} (Nahuel Defossé)\\\midrule
     15:40 \textendash\ 16:20 & Break\\\midrule
     16:20 \textendash\ 17:00 & Concluding discussion\\\midrule
     17:00 \textendash\ 19:00 & Happy Hour\\\midrule
     20:00 \textendash\ 23:00 & Dinner\\
     \bottomrule
   \end{tabular}
\end{center}

\clearpage

\subsection*{Locations}
% TODO Need to add more information
\begin{description}
   \item[Room 1] Sala de Projeção Harry Laus
   \item[Room 2] Sala de Projeção Henrique da Silva Fontes
   \item[Auditorium] Auditório Elke Hering
   \item[CIASC] Rua Murilo Andriani, 327 - Bairro Itacorubi. Florianópolis - SC
\end{description}

%\clearpage
% \section*{Program Details}
%
% \subsection*{Tutorials}
% \begin{description}
%    \item[16/05/2016] \ 
%    \begin{itemize}
%       \item \textbf{Python for geoscientist}, \emph{Filipe Fernandes (ocefpaf@gmail.com)}. Intermediate. Sala de projeção Harry Laus, 9h-12h.
%       \item \textbf{Utilización de Vispy para visualización rápida en 3D}, \emph{David Ochoa (ochoadavid@gmail.com)}. Intermediate. Sala de Projeção Harry Laus, 14h-17h.
%       \item \textbf{Introduction to MPI on python using mpi4py}, \emph{Jeudy Blanco (jeudyx@gmail.com)}. Introductory. Sala de Projeção Henrique da Silva Fontes, 9h-12h.
%       \item \textbf{Introduction to High-Performance Scientific Computing with PyOpenCL}, \emph{Celia Cintas (cintas.celia@gmail.com)}. Introductory. Sala de Projeção Henrique da Silva Fontes, 14h-17h.
%       \item \textbf{I see 42 everywhere}, \emph{Fernando Masanori (fmasanori@gmail.com)}. Introductory. Auditório Elke Hering, 9h-12h.
%       \item \textbf{The stupid content tracker}, \emph{Tomas Aliaga (tomas.aliaga@gmail.com)}. Introductory. Auditório Elke Hering, 14h-17h.
%    \end{itemize}
%    \item[17/05/2016] \ 
%    \begin{itemize}
%       \item \textbf{Fit and predict your data: Introduction to Scikit-learn}, \emph{Celia Cintas (cintas.celia@gmail.com)}. Introductory. Sala de Projeção Harry Laus, 9h-12h.
%       \item \textbf{scikit-image: Realce e Detecção de bordas}, \emph{Darleison Rodrigues (darleison.f@gmail.com)}. Intermediate. Sala de Projeção Harry Laus, 14h-17h. \underline{(Waiting for confirmation)}
%       \item \textbf{Scientific Python applied to Petroleum Engineering - an Introduction}, \emph{Igor T. Ghisi (igortg.subs@gmail.com)}. Introductory. Sala de Projeção Henrique da Silva Fontes, 9h-12h. \underline{(Waiting for confirmation)} 
%       \item \textbf{Time Series Analysis Using Python}, \emph{Bargava Subramanian (bargava@gmail.com)}. Introductory. Sala de Projeção Henrique da Silva Fontes, 14h-17h. \underline{(Waiting for confirmation)}
%       \item \textbf{Introduction to Data Analysis using Python Data Stack}, \emph{Bargava Subramanian (bargava@gmail.com)}. Introductory. Auditório Elke Hering, 9h-12h. \underline{(Waiting for confirmation)}
%     \item \textbf{Usando R com o Python - rpy2}, \emph{Arnaldo Russo (arnaldorusso@gmail.com)}. Intermediate. Auditório Elke Hering, 14h-17h. \underline{(Waiting for confirmation)}
%    \end{itemize}
% \end{description}

% Contra-capa
\newgeometry{left=0cm, bottom=0cm, top=0cm, right=0cm}
\thispagestyle{empty}
\noindent  % To remove the unwanted white space.
\includegraphics{contra-capa}
\NoBgThispage

\end{document}
